\documentclass[spanish,a4paper]{report}

\usepackage{polyglossia}
\usepackage{fontspec}

\setdefaultlanguage{spanish}
\setmainfont{CMU Serif}
\setsansfont{CMU Sans Serif}
\setmonofont{CMU Typewriter Text}

\usepackage{ccicons}
\usepackage{amsmath}
\usepackage{amssymb}
\usepackage{amsthm}
\usepackage{tcolorbox}
\tcbuselibrary{breakable}
\usepackage{cancel}

\usepackage{pgfplots}
\usepackage{tikz}
% this line is recomended by the pgfplots lastes stable version
% see https://pgfplots.sourceforge.net/pgfplots.pdf
\pgfplotsset{compat=1.18}

\usetikzlibrary{arrows.meta}

% This file keeps the main.tex file clean. Here is the declaration 
% of every environment I use in the document

% Config of the \newtheorem envs
% This avoids the italic font in theorem
% envs
\theoremstyle{definition}

% Definition environment. Follows the structure
% Definición <chapter><section><number>
\newtheorem{definition}{Definición}[section]

% A subtle teal frame, preserving theorem titles, numbering and references.
\tcolorboxenvironment{definition}{
  breakable,
  colback=teal!3!white,
  colframe=teal!55!white,
  boxrule=0.5pt,
  arc=2pt,
  left=10pt,
  right=10pt,
  top=6pt,
  bottom=6pt,
  before skip=12pt,
  after skip=12pt,
}

% Property environment. Follows the structure
% Propiedad <chapter><section><number>
\newtheorem{property}{Propiedad}[section]

% Theorem environment
% Teorema <chapter><section><number>
\newtheorem{theorem}{Teorema}[section]

% Example environment. Follows the structure
% Ejemplo <chapter><section><A,B,C,...> 
% A,B,C requieres a new counter.
\theoremstyle{remark}
\newtheorem{example}{Ejemplo}[section]
\renewcommand{\theexample}{\thesection.\Alph{example}}

% Redefinir el cierre para agregar el símbolo de fin
\let\origendexample\endexample
\def\endexample{\hfill\diamond\origendexample}

% Rename proof to Demostracion
\addto\captionsspanish{\renewcommand{\proofname}{Demostración}}

% Same as example, but with exercise
\newtheorem{exercise}{Ejercicio}[section]
\renewcommand{\theexercise}{\thesection.\Alph{exercise}}


\title{Mecánica de Órbitas para simuladores}
\author{Elio Valentino Anci}

\begin{document}
  \maketitle

  \begin{abstract}
    Este documento cubre las matemáticas necesarias para calcular y ejecutar distintas maniobras orbitales, incluyendo:
    \begin{enumerate}
      \item Cálculo del apoapsis (AP) y periapsis (PE) de una órbita.
      \item Modificación del apoapsis y periapsis mediante maniobras orbitales.
      \item Cambios de plano orbital para alinear la órbita con la de un objetivo.
      \item Maniobras de fase para sincronizar el encuentro entre dos naves.
      \item Transferencias interplanetarias.
    \end{enumerate}
  \end{abstract}

  \tableofcontents
  
  \chapter{Herramientas matemáticas}

En este capítulo se introducirán las matemáticas y geometría necesaria para comprender el resto de capítulos. Se verá una breve introducción a la geometría analítica y a vectores.

\section{Vectores}

En primer lugar se desarrollan conceptos relativos a \emph{espacios vectoriales}, estudiando sus características y propiedades. Esto será muy útil ya que proporcionará una estructura matemática para describir magnitudes como velocidad, posición, aceleración, entre otras; que serán especialmente útiles en la mayor parte del libro.

A continuación se presentarán geométrica y analíticamente los vectores en espacios bi y tridimensional, donde se dan las definiciones y las operaciones entre ellos, estableciendo las propiedades básicas.

\subsection{Espacios Vectoriales}

Para poder definir un espacio vectorial primero debe conocer la definición de vector,
\begin{definition}[Vector]
  Un vector \(\vec{v}\) en el plano xy es un par ordenado de números reales \((v_1,v_2)\). Los números \(v_1\) y \(v_2\) reciben el nombre de componentes del vector \(\vec{v}\).
\end{definition}



  \chapter{Glosario y vocabulario}\label{chpt:glosario}

En esta sección se explican, a grandes rasgos, algunos de los términos utilizados a lo largo del documento. El glosario se ha colocado al final para no interrumpir la lectura; cuando se introduzca vocabulario poco común, se podrá hacer referencia a esta sección para consultar su significado.

\section{Apoapsis y periapsis}

\paragraph{Apoapsis}
El apoapsis (también llamado \emph{apoastro}) es el punto de una órbita elíptica en el que el objeto orbitante se encuentra a mayor distancia del cuerpo central. Este punto puede recibir distintos nombres dependiendo del cuerpo alrededor del cual se orbita. Por ejemplo, si el cuerpo central es la Tierra se denomina \emph{apogeo}, mientras que si es el Sol se denomina \emph{afelio}.

En la figura \ref{fig_apoapsis} se representa el apoapsis de una órbita con una excentricidad relativamente grande. La excentricidad es una medida de cuánto se aparta la forma de una órbita de una circunferencia. 

\begin{figure}[!ht]
  \centering
  \begin{tikzpicture}[>=Stealth]
    \draw[fill=orange] (0,0) circle (6pt);
    \draw[dashed,thick] (0.92,0) ellipse (2.2 and 2);
    \draw[fill=blue!20] (3.12,0) circle (3pt);
    \draw[<-] (3.3,0) -- ++(.6,0) node[right]{Apoapsis};
  \end{tikzpicture}
  \caption{El cuerpo central (círculo anaranjado) se encuentra en uno de los
  focos de la elipse. El apoapsis corresponde al punto de la órbita en el que
  el objeto orbitante (círculo azul) se encuentra a mayor distancia del cuerpo
  central.}
  \label{fig_apoapsis}
\end{figure}

\paragraph{Periapsis}
El periapsis (también llamado \emph{periastro}) es, de manera análoga al apoapsis, el punto de una órbita elíptica en el que el objeto orbitante se encuentra a menor distancia del cuerpo central. Al igual que el apoapsis, puede recibir distintos nombres dependiendo del cuerpo central. Si este es la Tierra se denomina \emph{perigeo}, mientras que si es el Sol se denomina \emph{perihelio}.

\begin{figure}[!ht]
  \centering
  \begin{tikzpicture}[>=Stealth]
    \draw[fill=orange] (0,0) circle (6pt);
    \draw[dashed,thick] (0.92,0) ellipse (2.2 and 2);
    \draw[fill=blue!20] (-1.28,0) circle (3pt);
    \draw[<-] (-1.48,0) -- ++(-.6,0) node[left]{Periapsis};
  \end{tikzpicture}
  \caption{El periapsis corresponde al punto de la órbita en el que el objeto
  orbitante se encuentra a menor distancia del cuerpo central.}
  \label{fig_periapsis}
\end{figure}

Nótese que el periapsis y el apoapsis corresponden a puntos específicos de una órbita elíptica, y que el cuerpo central se encuentra situado en uno de los focos de la elipse.

El caso de una órbita perfectamente circular es especial. En ella todos los puntos de la órbita se encuentran a la misma distancia del cuerpo central, por lo que no existe un punto único que pueda distinguirse como apoapsis o periapsis. En este caso, las distancias correspondientes al apoapsis y al periapsis son iguales.


\end{document}
